% =====================================================================
%  Bourgain-pedia digestion of:
%    J. Bourgain, L. Clozel, J.-P. Kahane,
%    "Principe d'Heisenberg et fonctions positives",
%    Ann. Inst. Fourier 60, No. 4, 1215-1232 (2010).
%    arXiv:0811.4360 -- DOI 10.5802/aif.2552 -- Zbl 1298.11105
%
%  Original: 19 pages (the authors' own LaTeX, kept byte-exact in
%  ...-original.tex).  Ledger: ledger.md.  Notes: notes.md.
% =====================================================================
\documentclass[11pt,reqno]{amsart}
\usepackage[T1]{fontenc}
\usepackage{amsmath,amssymb,amsthm}
\usepackage[margin=1.05in]{geometry}
\usepackage{xcolor}
\usepackage[most]{tcolorbox}
\usepackage{enumitem}
\usepackage{hyperref}
\hypersetup{colorlinks=true,linkcolor=black,urlcolor=teal}

% ----- the four editorial environments -------------------------------
\definecolor{expcol}{HTML}{2F5D7C}
\definecolor{filcol}{HTML}{2E6B4F}
\definecolor{notcol}{HTML}{6B4E9B}
\definecolor{caucol}{HTML}{9C3B22}

\newtcolorbox{expansion}{breakable,enhanced,colback=expcol!4,colframe=expcol!55,
  boxrule=0.35pt,left=6pt,right=6pt,top=4pt,bottom=4pt,
  fonttitle=\bfseries\scriptsize,title={EXPANSION \textnormal{--- a step they compressed, written out}}}
\newtcolorbox{filled}{breakable,enhanced,colback=filcol!4,colframe=filcol!55,
  boxrule=0.35pt,left=6pt,right=6pt,top=4pt,bottom=4pt,
  fonttitle=\bfseries\scriptsize,title={FILLED \textnormal{--- an argument omitted in the original, supplied here}}}
\newtcolorbox{ournotation}{breakable,enhanced,colback=notcol!4,colframe=notcol!55,
  boxrule=0.35pt,left=6pt,right=6pt,top=4pt,bottom=4pt,
  fonttitle=\bfseries\scriptsize,title={OUR NOTATION \textnormal{--- not in the original}}}
\newtcolorbox{caution}{breakable,enhanced,colback=caucol!5,colframe=caucol!70,
  boxrule=0.8pt,left=6pt,right=6pt,top=4pt,bottom=4pt,
  fonttitle=\bfseries\scriptsize,title={CAUTION \textnormal{--- imported, or not fully justified here}}}

\theoremstyle{plain}
\newtheorem*{thmA}{Theorem 1}
\newtheorem*{thmB}{Theorem 2}
\newtheorem*{thmC}{Theorem 3}
\newtheorem*{thmD}{Theorem 4}
\newtheorem*{lemA}{Lemma 1}
\newtheorem*{lemB}{Lemma 2}
\newtheorem*{propA}{Proposition 1}
\newtheorem*{propB}{Proposition 2}
\newtheorem*{propOurs}{Proposition (ours)}
\theoremstyle{definition}
\newtheorem*{bg}{Background}
\theoremstyle{remark}
\newtheorem*{rmk}{Remark}

\newcommand{\R}{\mathbb{R}}
\newcommand{\Z}{\mathbb{Z}}
\newcommand{\Q}{\mathbb{Q}}
\newcommand{\C}{\mathbb{C}}
\newcommand{\Ss}{\mathcal{S}}
\newcommand{\Bcal}{\mathcal{B}}
\newcommand{\fh}{\widehat{f}}
\DeclareMathOperator{\supp}{supp}
\DeclareMathOperator{\Res}{Res}
\DeclareMathOperator{\Tr}{Tr}
\newcommand{\Ledger}[1]{{\scriptsize\textsf{[#1]}}}

\begin{document}

\title[Heisenberg principle and positive functions --- digestion]
      {Principe d'Heisenberg et fonctions positives\\
       \normalsize The Heisenberg principle and positive functions\\
       \normalsize a Bourgain-pedia digestion}
\author{J. Bourgain, L. Clozel and J.-P. Kahane}
\address{Original: Ann.\ Inst.\ Fourier \textbf{60} (2010), 1215--1232; arXiv:0811.4360}
\date{Digestion generated \today}
\maketitle

% =====================================================================
\section*{\S0. Editorial preface}
% =====================================================================

This is the Bourgain-pedia \emph{digestion} of Bourgain, Clozel and Kahane's
\emph{Principe d'Heisenberg et fonctions positives}: the paper at full length, with
every step they compressed written out and every argument they omitted supplied.

\medskip
\noindent\textbf{Source.} Unusually for this project, the source is the authors' own
LaTeX, retrieved from the arXiv e-print service and kept byte-exact as
\texttt{...-original.tex}. Nothing here is a transcription, so no formula below can
have been corrupted by us in the reading; the slips listed are demonstrably theirs.
The original is \textbf{19 pages}.

\medskip
\noindent\textbf{Language.} The paper is in French; this digestion is in English,
because the project is. Quotations are given in French with a translation. This is an
editorial choice, not a rule of the blueprint --- see \texttt{notes.md}.

\medskip
\noindent\textbf{How to read the boxes.} Four kinds of material are ours, and each is
boxed: \emph{expansion} (a step of theirs, written out), \emph{filled} (an argument
absent from the original), \emph{our notation}, and \emph{caution} (something imported
that we have not verified). Unboxed prose and displays are theirs.

\medskip
\noindent\textbf{Everything numerical here has been recomputed independently}, before
this text was drafted: $\lambda$; both constants of Theorem 1; every coefficient of
both Taylor expansions; $P_4(3/2)=3/2$; $Q_4(c+1)=-c^2+1$; the $a=\sqrt2$ and $a=2$
reductions; $p_5=0$ and $p_6=-4/45$; and Theorem 3's inequality chain in nine
dimensions. All agree with the paper.

\medskip
\noindent\textbf{Seven slips in the original}, none affecting a theorem:
\begin{enumerate}[label=(\roman*),leftmargin=2.2em]
\item \S1: ``on en d\'eduit que $\Bcal_1^-\le B_1$'' where $\Bcal_1^-\le B_1^-$ is
  needed. \Ledger{B20}
\item \S1: ``$B_1^-\le(a\sqrt2)^2=2a$'' --- the square is $2a^2$. \Ledger{B27}
\item \S2: $e^{X(2h-3h^2+3h^3-4h^4)X}$ carries a doubled $X$; and the series is
  written $P_1h+P_2h^2+P_3h^3+P_4h^5$, where $P_4h^4$ is meant. \Ledger{C11}
\item \S2: the extremum condition is printed $a(1-a)=2\log(1+a)$, which has no root
  $>1$; it must be $a(a-1)=2\log(1+a)$, whose root \emph{is} the $2.08137\ldots$
  printed. \Ledger{C18}
\item \S3: the label $(3.4)$ is used twice, for two different equations.
  \Ledger{D5}
\item \S4: Odlyzko's bound is written $|D|^{1/d}\ge22.2(1+0(d))$; $o(1)$ is meant.
  \Ledger{E21}
\item \S4: $(4.9)$ typesets $C=|D|^{-1_/d_0}$; and the sentence after it cites
  ``(3.10) et (4.8)'' where $(4.9)$ is meant. \Ledger{E25, E26}
\end{enumerate}

\smallskip
\noindent Two places where a stated bound is weaker than the paper's own method
delivers --- both flagged by the authors themselves as approximate, and both recorded
in \S5 rather than silently improved:
\begin{itemize}[leftmargin=1.6em]
\item \S2: ``une combinaison lin\'eaire convenable de $H_0$ et $H_4$ donne
  $\pi A^2\le 3$''. Carried out (\S2.1, which the paper does not do), the combination
  gives $\pi A^2=3/2$ --- exactly the bound $(2.5)$ obtained later by a different
  route. \Ledger{C2}
\item \S2: the correction argument is stated to give $A\le0.63\ldots$, ``nous n'avons
  fait qu'un calcul tr\`es approch\'e''. Optimised numerically it gives
  $A\le0.5973$. \Ledger{C21}
\end{itemize}

\medskip
\noindent\textbf{What we do not verify.} All of \S4 rests on Tate's thesis and on four
further citations (Armitage, Serre, Odlyzko, Roquette/Golod--Shafarevi\v c). We state
each import precisely and check the paper's use of it, but we have not read those
sources; the \texttt{caution} boxes in \S4 say so at each point.

% =====================================================================
\section*{\S0$'$. Notation}
% =====================================================================

\noindent\textbf{Theirs.}
\begin{itemize}[leftmargin=1.6em,itemsep=1pt]
\item $\fh(y)=\int f(x)e^{-2i\pi xy}\,dx$, with inverse $f(x)=\int\fh(y)e^{2i\pi xy}dy$;
  in $\R^d$, $xy$ becomes $x\cdot y$ and $dx$ Lebesgue measure. This convention makes
  $\gamma(x)=e^{-\pi x^2}$ (resp.\ $e^{-\pi\|x\|^2}$) self-dual, and every constant
  below depends on it. \Ledger{A6}
\item A \emph{Fourier pair} $(f,\fh)$: both in $L^1$, so both continuous and vanishing
  at infinity.
\item $a_f,a_{\fh}$: any thresholds beyond which $f$, resp.\ $\fh$, is $\ge0$.
\item $A(f)=\inf\{x>0: f\ge0 \text{ on } (x,\infty)\}$ --- the least such threshold;
  in $\R^d$, $\inf\{r>0: f(x)\ge0 \text{ for } \|x\|>r\}$.
\item $f^\#$: the spherical average of $f$ (\S3).
\end{itemize}

\begin{ournotation}
\textbf{The four constants}, which the paper introduces over two pages and never
tabulates. All are infima over Fourier pairs satisfying (1--3); $\Ss$ is the Schwartz
class. \Ledger{A2}
\[
\begin{array}{lll}
 B_1 &=\inf A(f)A(\fh) & \text{over } f\in L^1\\
 \Bcal_1 &=\inf A(f)A(\fh) & \text{over } f\in\Ss\\
 B_1^- &=\inf A(f)A(\fh) & \text{over } f\in L^1 \text{ with } f(0)<0\\
 \Bcal_1^- &=\inf A(f)A(\fh) & \text{over } f\in\Ss \text{ with } f(0)<0
\end{array}
\]
and $B_d,\Bcal_d$ likewise in $\R^d$. \S1 proves $B_1\le\Bcal_1\le2B_1$, so the two
that matter differ by at most a factor 2. Note that in the source \verb|\bb| and
\verb|\Bb| are both defined as $\mathcal B$, so $\Bcal_d$ appears under two spellings
that typeset identically. \Ledger{F1}
\end{ournotation}

\begin{ournotation}
We write $f^+=\max(f,0)$, $f^-=\max(-f,0)$, so $f=f^+-f^-$ and $|f|=f^++f^-$; the
paper uses this without introducing it. We reserve $\gamma$ for the Gaussian in
whatever dimension is current. \Ledger{F5}
\end{ournotation}

% =====================================================================
\section*{\S0$''$. Background}
% =====================================================================

\begin{bg}[Inversion and the elementary bounds]\label{bg:inv}
If $f,\fh\in L^1(\R^d)$ then both are continuous and tend to $0$ at infinity, and
$\|f\|_\infty\le\|\fh\|_1$, $\|\fh\|_\infty\le\|f\|_1$. In particular
$\fh(0)=\int f$ and $f(0)=\int\fh$.
\end{bg}

\begin{bg}[Dilation]\label{bg:dil}
For $\lambda>0$, if $g(x)=f(x/\lambda)$ then $\widehat g(y)=\lambda^d\fh(\lambda y)$.
Hence $A(g)=\lambda A(f)$ and $A(\widehat g)=\lambda^{-1}A(\fh)$, so the product
$A(f)A(\fh)$ is dilation-invariant.
\end{bg}

\begin{bg}[Hermite functions \Ledger{C1}]\label{bg:herm}
There is an orthogonal basis $(H_n)_{n\ge0}$ of $L^2(\R)$ with
$H_n(x)=e^{-\pi x^2}P_n(x)$, $\deg P_n=n$, and $\widehat{H_n}=i^{-n}H_n$ for the
convention above. Consequently a self-dual $f=\fh$ has an expansion supported on
$n\equiv0\pmod4$.
\end{bg}

\begin{bg}[A function and its transform cannot both have compact support]\label{bg:cpt}
If $f\in L^1(\R^d)$ and both $f$ and $\fh$ have compact support then $f=0$. (The
transform of a compactly supported $L^1$ function extends to an entire function on
$\C^d$; an entire function vanishing on a non-empty open set vanishes identically.)
\end{bg}

\begin{bg}[Stirling, lower bound]\label{bg:stir}
For $m>0$, $\Gamma(m+1)>\sqrt{2\pi m}\,(m/e)^m$.
\end{bg}

% =====================================================================
\section*{\S1. The problem, and a lower bound for $B_1$}
% =====================================================================

\begin{ournotation}
\textbf{The paper's own summary}, in place of an abstract, translated. \Ledger{A1}
``We describe a natural problem concerning Fourier transforms, which introduces
constants $B_d$ depending on the dimension $d$. The problem is posed, in dimension 1,
in the first section, where we describe a simple reduction to the case of self-dual
functions. The first result is a lower bound for $B_1$. Then we consider a very
natural class of functions which gives a simple upper bound for $B_1$. This is not
optimal: we describe simple arguments which allow it to be improved (\S2). In \S3
these calculations are extended to arbitrary dimensions, giving simple lower and upper
bounds for the constants. Finally \S4, arithmetic, relates this problem to a well-known
question concerning zeta functions of number fields. The arithmetic arguments show that
the linear growth of $B_d$ as a function of the dimension is natural in view of known
properties of the ramification of these fields.''
\end{ournotation}

\subsection*{1.1 The conditions}

Consider a Fourier pair $(f,\fh)$ on the real line, and the conditions
\begin{enumerate}[label=\textbf{\arabic*)},leftmargin=2.4em]
\item $f$ and $\fh$ are real and even, not identically zero;
\item $f(0)\le0$ and $\fh(0)\le0$;
\item $f(x)\ge0$ for $x\ge a_f$, and $\fh(y)\ge0$ for $y\ge a_{\fh}$.
\end{enumerate}

\noindent\textbf{Problem.} What is the infimum of the product $a_fa_{\fh}$ over
Fourier pairs satisfying (1--3)? Call it $B_1\ge0$.

\begin{filled}
\textbf{Why $a_f,a_{\fh}>0$.} \Ledger{A3} The paper asserts that 2) and $f\not\equiv0$
force $a_f,a_{\fh}>0$, without argument. Suppose $a_f=0$ were admissible, i.e.\
$f\ge0$ on all of $(0,\infty)$; $f$ is even, so $f\ge0$ everywhere except possibly at
$0$, and by continuity $f\ge0$ everywhere. Then $\fh(0)=\int f\ge0$, and 2) gives
$\fh(0)\le0$, so $\int f=0$ with $f\ge0$ continuous, whence $f\equiv0$ --- excluded by
1). The same argument applies to $\fh$.
\end{filled}

\begin{filled}
\textbf{Such pairs exist, and the problem is not vacuous.} \Ledger{A4} The paper says
only ``de tels couples existent \`a l'\'evidence''. Here is one. Take
$f(x)=x^2\bigl(x^2-\tfrac{3}{2\pi}\bigr)e^{-\pi x^2}$. \S2.1 below verifies that
$f=\fh$, that $f(0)=0$, and that $f\ge0$ exactly where $x^2\ge3/(2\pi)$. So $(f,f)$
satisfies (1--3) with $a_f=a_{\fh}=\sqrt{3/(2\pi)}$, giving at once
$B_1\le3/(2\pi)=0.477\ldots$

\textbf{And why positivity is not obvious.} \Ledger{A5} The pair may be dilated:
$f(x/\lambda)$ has transform $\lambda\fh(\lambda y)$, which sends
$(a_f,a_{\fh})\mapsto(\lambda a_f,\lambda^{-1}a_{\fh})$. So each threshold separately
can be made as small as one likes, and only the \emph{product} is constrained.
Nothing elementary prevents the product from being driven to $0$ as well; that it
cannot be is Theorem 1.
\end{filled}

\subsection*{1.2 Reduction to self-dual functions with $f(0)=0$}

The product $A(f)A(\fh)$ is invariant under $f(x)\mapsto f(x/\lambda)$,
$\fh(y)\mapsto\lambda\fh(\lambda y)$ (Background \ref{bg:dil}). Since
$B_1=\inf A(f)A(\fh)$, one may restrict to pairs with $A(f)=A(\fh)$. Then
$f+\fh\ne0$, and
\[A(f+\fh)\le A(f)=A(\fh),\]
so $B_1=\inf A^2(f+\fh)$, and therefore
\begin{equation}\label{eq:BA}
 B_1=A^2,\qquad A=\inf A(f),
\end{equation}
the infimum over $f\in L^1(\R)$ real, even, not identically zero, equal to its own
Fourier transform, with $f(0)\le0$.

\begin{expansion}
\textbf{The three steps, done.} \Ledger{B1, B2, B4, B5}

\emph{(a) Normalising $A(f)=A(\fh)$.} Given a pair, put
$\lambda=\bigl(A(\fh)/A(f)\bigr)^{1/2}$ and replace $f(x)$ by $f(x/\lambda)$. By
Background \ref{bg:dil} the new thresholds are $\lambda A(f)$ and
$\lambda^{-1}A(\fh)$, and both equal $\sqrt{A(f)A(\fh)}$. The product is unchanged,
so the infimum is unchanged.

\emph{(b) $f+\fh$ is admissible, and $A(f+\fh)\le A(f)$.} It is real and even;
$(f+\fh)^{\wedge}=\fh+f$ because $\widehat{\fh}=f$ for even $f$; and
$(f+\fh)(0)=f(0)+\fh(0)\le0$. For $x>A(f)=A(\fh)$ both summands are $\ge0$, so the
sum is, giving $A(f+\fh)\le A(f)$.

\emph{(c) $f+\fh\not\equiv0$.} The paper's parenthesis is ``consider its values at
points near $A(f)$ and above it''. \Ledger{B3} In full: by definition of $A(f)$ as an
infimum, $f$ takes a strictly negative value at some point arbitrarily close to
$A(f)$ from below --- otherwise $f\ge0$ on a larger interval and $A(f)$ was not the
infimum. Take $x_0<A(f)$ with $f(x_0)<0$. If $f+\fh$ vanished identically then
$\fh=-f$, so $\fh(x_0)=-f(x_0)>0$; but also $\fh=-f$ gives
$f=\widehat{\fh}=-\fh=f$, no contradiction yet. Instead use positivity: for
$x>A(f)$, $f(x)\ge0$ and $\fh(x)\ge0$, and $f+\fh\equiv0$ forces both to vanish on
$(A(f),\infty)$. Then $f$ vanishes on a half-line, and being even, on $|x|>A(f)$:
$f$ has compact support. The same holds for $\fh$, so $f\equiv0$ by Background
\ref{bg:cpt} --- contradicting 1).

\emph{(d) Both inequalities for $B_1=\inf A^2(f+\fh)$.} \Ledger{B5} ``$\le$'': each
$f+\fh$ arising from an admissible pair is itself admissible and self-dual, so
$A^2(f+\fh)=A(f+\fh)A(\widehat{f+\fh})$ is one of the products over which $B_1$ is an
infimum. ``$\ge$'': by (b), $A^2(f+\fh)\le A(f)A(\fh)$ for every normalised pair, so
the infimum over the self-dual ones is at most $B_1$.
\end{expansion}

Setting $\gamma(x)=e^{-\pi x^2}$, so that $\gamma=\widehat\gamma$: if $f(0)<0$ then
$f-f(0)\gamma$ is not identically zero, satisfies the same conditions, and
$A(f-f(0)\gamma)\le A(f)$. Finally
\begin{equation}\label{eq:11}
 A=\inf A(f),
\end{equation}
\emph{the infimum being over $f\in L^1(\R)$ real, even, not identically zero, with
$f=\fh$ and $f(0)=0$.}

\begin{expansion}
\textbf{The Gaussian correction.} \Ledger{B6} Since $f(0)<0$, $-f(0)>0$, so
$f-f(0)\gamma=f+|f(0)|\gamma$ \emph{adds} a positive multiple of a positive function.
Hence: it is self-dual (both summands are); its value at $0$ is
$f(0)-f(0)\gamma(0)=f(0)-f(0)=0$; it is $\ge f$ pointwise, so it is $\ge0$ wherever
$f$ is, giving $A(f-f(0)\gamma)\le A(f)$; and it is not identically zero, since
otherwise $f=f(0)\gamma$ would be everywhere $\le0$, contradicting 3) as in the
argument for $a_f>0$.
\end{expansion}

\subsection*{1.3 Theorem 1}

\begin{thmA}
Let $\lambda=-\inf\bigl(\tfrac{\sin x}{x}\bigr)=0{,}2172\cdots$. Then
\[A\ge\frac{1}{2(1+\lambda)}=0{,}4107\cdots,\qquad\text{hence}\qquad
  B_1\ge0{,}1687\cdots\]
\end{thmA}

\begin{expansion}
\textbf{The constant.} \Ledger{B8, B9} $\lambda$ is attained at the first positive
critical point of $\sin x/x$, i.e.\ the first positive root of $\tan x=x$.
Recomputed: $x=4{,}4934094579$, $\lambda=0{,}2172336282$, and then
$1/(2(1+\lambda))=0{,}4107674882$ and $A^2\ge0{,}1687299294$. Both agree with the
printed values.
\end{expansion}

\begin{proof}
Choose $f=\fh$ with $f(0)=0$ and $\int_\R|f|=1$; write $A=A(f)$. Put $f=f^+-f^-$,
$|f|=f^++f^-$. Since $\int_\R f=\fh(0)=0$ we have
$\int f^+=\int f^-=\int_{-A}^{A}f^-=\tfrac12$. Hence
$\int_{|x|\ge A}|f|=\int_{|x|\ge A}f^+\le\tfrac12$, so
$\int_{|x|\le A}|f|\ge\tfrac12$. But $|f(x)|\le\int|\fh|=1$. Therefore
$2A\ge\tfrac12$, a first bound $A\ge\tfrac14$. This argument extends to higher
dimensions (\S3).

In dimension 1 it can be refined as follows. Since $f$ is self-dual,
\[
\begin{aligned}
 f(x)&=\int f(y)\cos2\pi yx\,dy=\int f(y)(\cos2\pi yx-1)\,dy\\
     &=\int f^-(y)(1-\cos2\pi yx)\,dy-\int f^+(y)(1-\cos2\pi yx)\,dy .
\end{aligned}
\]
Both integrals being positive, this implies
\[f^-(x)\le\int f^+(y)(1-\cos2\pi yx)\,dy,\]
whence
\[\tfrac14=\int_0^Af^-\le\int_{-\infty}^{+\infty}f^+(y)\Bigl[A-\frac{\sin2\pi yA}{2\pi y}\Bigr]dy\]
and therefore
\[\tfrac14\le\frac A2\sup_{u\in\R}\Bigl(1-\frac{\sin u}{u}\Bigr)=\frac A2(1+\lambda),\]
which is the theorem.
\end{proof}

\begin{expansion}
\textbf{Every step of that proof.} \Ledger{B10--B18}

\emph{Normalising $\int|f|=1$} is legitimate: $A(cf)=A(f)$ for $c>0$, and
$f\mapsto cf$ preserves self-duality and $f(0)=0$. \Ledger{B10}

\emph{$\int f^+=\int f^-=\tfrac12$.} \Ledger{B11} From $\int f=\fh(0)=0$ we get
$\int f^+=\int f^-$, and $\int f^++\int f^-=\int|f|=1$. That $\int f^-$ is
concentrated on $[-A,A]$ is because $f\ge0$ outside $[-A,A]$ (evenness plus the
definition of $A$), so $f^-$ vanishes there.

\emph{$\int_{|x|\le A}|f|\ge\tfrac12$.} \Ledger{B12} On $|x|\ge A$ we have $f\ge0$, so
$|f|=f^+$ there, and $\int_{|x|\ge A}f^+\le\int f^+=\tfrac12$; subtract from
$\int|f|=1$.

\emph{$\|f\|_\infty\le1$ and $A\ge\tfrac14$.} \Ledger{B13} By Background \ref{bg:inv},
$\|f\|_\infty\le\|\fh\|_1=\|f\|_1=1$ using $f=\fh$. Then
$\tfrac12\le\int_{|x|\le A}|f|\le 2A\cdot1$.

\emph{The cosine identity.} \Ledger{B15} $f$ real and even gives
$\fh(x)=\int f(y)\cos2\pi yx\,dy$; self-duality replaces $\fh$ by $f$ on the left.
Subtracting $\int f(y)\,dy=0$ inserts the $-1$.

\emph{The inequality for $f^-$.} \Ledger{B16} Write $I^\pm(x)=\int
f^\pm(y)(1-\cos2\pi yx)dy$, both $\ge0$ since $1-\cos\ge0$ and $f^\pm\ge0$. The
identity is $f=I^--I^+$, so $-f=I^+-I^-\le I^+$. Also $0\le I^+$. Hence
$f^-=\max(-f,0)\le I^+$, which is the display.

\emph{The Fubini step.} \Ledger{B17} $\int_0^Af^-=\tfrac14$ because $f^-$ is even with
total mass $\tfrac12$, supported in $[-A,A]$. Then
\[\int_0^A\!\!\int f^+(y)(1-\cos2\pi yx)\,dy\,dx
 =\int f^+(y)\Bigl[A-\frac{\sin2\pi yA}{2\pi y}\Bigr]dy,\]
the interchange being justified because the integrand is non-negative and
$f^+\in L^1$ (Tonelli), and the inner integral is
$\int_0^A(1-\cos2\pi yx)dx=A-\sin(2\pi yA)/(2\pi y)$.

\emph{The last line.} \Ledger{B18} Substituting $u=2\pi yA$,
$A-\sin(2\pi yA)/(2\pi y)=A\bigl(1-\frac{\sin u}{u}\bigr)\le A\sup_u(1-\frac{\sin
u}{u})=A(1+\lambda)$, and $\int f^+=\tfrac12$ gives the factor $A(1+\lambda)/2$.
\end{expansion}

\subsection*{1.4 The Schwartz class costs at most a factor 2}

Let $\Bcal_1=A^2$ with $A$ defined by \eqref{eq:11} but with $f\in\Ss$. Evidently
\begin{equation}\label{eq:12}
 B_1\le\Bcal_1 .
\end{equation}
Let $B_1^-=\inf(A^2\mid f(0)<0,\ f=\fh\text{ even}\ne0,\ f\in L^1)$, and $\Bcal_1^-$
likewise with $f\in\Ss$. Evidently
\begin{equation}\label{eq:13}
 B_1^-\le\Bcal_1^-,
\end{equation}
\begin{equation}\label{eq:14}
 \Bcal_1\le\Bcal_1^-,\qquad B_1\le B_1^- .
\end{equation}

Let us check that $\Bcal_1^-\le B_1^-$. Let $f\in L^1$ satisfy the condition of
\eqref{eq:11} but with $f(0)<0$, and let $a=A(f)$. Let $\varphi=\psi*\psi$ with $\psi$
smooth, even, positive, of compact support very close to $0$, and $g=f*\varphi$. Then
$A(g)\le a+\varepsilon$ and $g(0)<0$. We have $\widehat g=\fh\widehat\psi^{\,2}$;
performing the same operation on $\widehat g$ we obtain $h\in\Ss$ with $h=\widehat h$,
$h(0)<0$ and $A(h)\le a+\varepsilon$; whence $\Bcal_1^-\le B_1^-$, that is
\begin{equation}\label{eq:15}
 \Bcal_1^-=B_1^- .
\end{equation}

\begin{expansion}
\textbf{Two things about that paragraph.} \Ledger{B20, B21}

\emph{The printed conclusion.} The original reads ``on en d\'eduit que
$\Bcal_1^-\le B_1$ soit $\Bcal_1^-=B_1^-$''. The intermediate inequality must be
$\Bcal_1^-\le B_1^-$: it is $B_1^-$, not $B_1$, that combines with \eqref{eq:13} to
give the stated equality. We have written $B_1^-$.

\emph{Why the operation is performed twice.} One convolution does not preserve
self-duality: $g=f*\varphi$ has $\widehat g=\fh\,\widehat\psi^{\,2}$, a
\emph{product}, not a convolution, so $g\ne\widehat g$ in general. Repeating on the
transform side --- multiplying $g$ by a suitable smooth positive even factor, which is
convolution on the transform side --- restores the symmetry. Concretely, with
$\psi$ real even one has $\widehat\psi$ real even, so
$h:=(f*\varphi)\cdot\widehat\varphi$ satisfies $\widehat h=\fh\widehat\varphi\cdot
\varphi=h$ using $f=\fh$ and $\varphi,\widehat\varphi$ real even. Both operations
move $A$ by at most $\varepsilon$ if $\supp\psi$ is small enough, and neither changes
the sign of the value at $0$. That $g(0)<0$ survives is because $\varphi\ge0$ has mass
$1$ and $f$ is continuous and negative near $0$.
\end{expansion}

Note that the argument does not apply if $f(0)=0$. We shall show
\begin{equation}\label{eq:16}
 B_1^-\le2\,B_1 ;
\end{equation}
by \eqref{eq:14} and \eqref{eq:16} we deduce
\begin{equation}\label{eq:17}
 B_1\le\Bcal_1\le2\,B_1 .
\end{equation}

\begin{expansion}
\textbf{The chain, checked.} \Ledger{B19, B28} The left inequality is \eqref{eq:12}.
For the right: $\Bcal_1\le\Bcal_1^-$ by \eqref{eq:14}, $\Bcal_1^-=B_1^-$ by
\eqref{eq:15}, and $B_1^-\le2B_1$ by \eqref{eq:16}; compose. Note that
\eqref{eq:15} is the only one of the four relations requiring work --- \eqref{eq:12},
\eqref{eq:13} and \eqref{eq:14} are immediate from the definitions, since each
enlarges the class over which the infimum is taken or drops a constraint.
\end{expansion}

Let $f$ satisfy \eqref{eq:11} and $a=A(f)$. Since $\fh(0)=\int f=0$, $f$ takes
strictly negative values on $[-a,a]$. Let $b>0$ with $f(b)<0$, and consider the
distribution
\[T=\delta_b+\delta_{-b}+2\delta_0 .\]
It is a positive measure, of positive type:
\[\widehat T=2\cos(2\pi by)+2\ge0 .\]
We have $(T*f)(0)=f(b)+f(-b)<0$. Since $b<a$, $g=T*f$ satisfies
\[g(0)<0,\qquad g\ge0\ \text{on}\ [2a,\infty),\]
and moreover $\widehat g=\widehat T\fh$ is $\ge0$ on $[a,\infty)$, with $\widehat
g(0)=0$. By dilation one then obtains $h$ with
\[h\ge0\ \text{on}\ [a\sqrt2,\infty),\quad h(0)<0;\qquad
  \widehat h\ge0\ \text{on}\ [a\sqrt2,\infty),\quad \widehat h(0)=0 .\]
The functions $h,\widehat h$ are real and even. Then $h+\widehat h$ satisfies the
conditions defining $B_1^-$. So $B_1^-\le(a\sqrt2)^2=2a^2$; varying $f$ gives
\eqref{eq:16}.

\begin{expansion}
\textbf{The convolution trick, in full.} \Ledger{B22--B27}

\emph{$T$ is of positive type.} $\widehat{\delta_c}(y)=e^{-2\pi icy}$, so
$\widehat T=e^{-2\pi iby}+e^{2\pi iby}+2=2\cos(2\pi by)+2\ge0$.

\emph{The value at $0$.} $(T*f)(0)=f(-b)+f(b)+2f(0)$. The printed line drops the last
term because $f(0)=0$ under \eqref{eq:11}; we say so rather than leaving the reader to
notice. \Ledger{B23}

\emph{$g\ge0$ on $[2a,\infty)$.} \Ledger{B24} $g(x)=f(x-b)+f(x+b)+2f(x)$. For
$x\ge2a$ and $0<b<a$, each of $x-b,x,x+b$ is $\ge2a-a=a$, so each term is $\ge0$.

\emph{The transform side.} $\widehat g=\widehat T\fh$ with $\widehat T\ge0$ and
$\fh\ge0$ on $[a,\infty)$, so $\widehat g\ge0$ there; and
$\widehat g(0)=\widehat T(0)\fh(0)=4\cdot0=0$.

\emph{The dilation.} \Ledger{B26} $g$ is non-negative from $2a$ and $\widehat g$ from
$a$: the two thresholds differ. Replace $g(x)$ by $g(x/\lambda)$, whose transform is
$\lambda\widehat g(\lambda y)$; the thresholds become $2a\lambda$ and $a/\lambda$.
They are equal when $\lambda=1/\sqrt2$, and their common value is then
$2a/\sqrt2=a\sqrt2$. This is the $h$ of the display. Note $h(0)=g(0)<0$ and
$\widehat h(0)=\lambda\widehat g(0)=0$.

\emph{Why $h+\widehat h$.} It is self-dual, even, real; $(h+\widehat h)(0)=h(0)<0$, so
it is a competitor for $B_1^-$; and both summands are $\ge0$ beyond $a\sqrt2$, so
$A(h+\widehat h)\le a\sqrt2$.

\emph{The printed value.} \Ledger{B27} $B_1^-\le(a\sqrt2)^2$ and $(a\sqrt2)^2=2a^2$;
the original prints $2a$. It is $2a^2$ that gives \eqref{eq:16}, since
$B_1=\inf a^2$ over the competitors of \eqref{eq:11}.
\end{expansion}

% =====================================================================
\section*{\S2. An upper bound for $B_1$}
% =====================================================================

\subsection*{2.1 The Hermite route, carried out}

A first idea is to associate to $f$ its Hermite expansion
$f\sim\sum_0^\infty a_nH_n(x)$, where the $H_n$ are eigenvectors of the Fourier
operator $\mathcal F$ with eigenvalues $i^n$; thus $f=\fh$ is expressed as
$f\sim\sum_0^\infty a_{4m}H_{4m}$. The $H_n$ have the form $H_n(x)=e^{-\pi x^2}P_n(x)$
with $P_n$ a polynomial of degree $n$. \emph{A suitable linear combination of $H_0$
and $H_4$ (such that $f(0)=0$) gives $\pi A^2\le3$.} Beyond this the calculations seem
difficult and we have not pursued this route.

\begin{filled}
\textbf{The combination, and what it actually gives.} \Ledger{C2} The paper states
$\pi A^2\le3$ and does not compute. Carrying it out is short, and the answer is
better than advertised.

Seek $f=(\alpha+\beta x^2+\delta x^4)\gamma$, $\gamma(x)=e^{-\pi x^2}$ --- the general
even element of the span of $H_0,H_2,H_4$. From
$\widehat{xf}=\frac{i}{2\pi}\,\widehat f{}'$ one gets
$\widehat{x^2\gamma}=-\frac1{4\pi^2}\gamma''$ and
$\widehat{x^4\gamma}=\frac1{16\pi^4}\gamma''''$, and differentiating
$\gamma'=-2\pi x\gamma$ repeatedly,
\[\gamma''=(4\pi^2x^2-2\pi)\gamma,\qquad
  \gamma''''=(16\pi^4x^4-48\pi^3x^2+12\pi^2)\gamma .\]
Hence
\[\widehat{x^2\gamma}=\Bigl(\tfrac1{2\pi}-x^2\Bigr)\gamma,\qquad
  \widehat{x^4\gamma}=\Bigl(x^4-\tfrac{3}{\pi}x^2+\tfrac{3}{4\pi^2}\Bigr)\gamma,\]
so that
\[\fh=\Bigl[\alpha+\tfrac{\beta}{2\pi}+\tfrac{3\delta}{4\pi^2}
   -\Bigl(\beta+\tfrac{3\delta}{\pi}\Bigr)x^2+\delta x^4\Bigr]\gamma .\]
Matching coefficients with $f$: the $x^4$ terms agree automatically; the $x^2$ terms
give $2\beta=-3\delta/\pi$; and the constant terms give
$\beta/(2\pi)+3\delta/(4\pi^2)=0$, which is the \emph{same} condition. So
self-duality is exactly $\beta=-\tfrac{3\delta}{2\pi}$, with $\alpha$ free, and
$f(0)=\alpha$, so the normalisation $f(0)=0$ forces $\alpha=0$. Taking $\delta=1$,
\[\boxed{\,f(x)=x^2\Bigl(x^2-\tfrac{3}{2\pi}\Bigr)e^{-\pi x^2}\,}\]
which is self-dual (verified numerically to 12 digits at five points), vanishes at
$0$, and is $\ge0$ exactly for $x^2\ge3/(2\pi)$. Hence
\[A\le\sqrt{\tfrac{3}{2\pi}}=0{,}6909883\ldots,\qquad \pi A^2=\tfrac32 .\]

So the Hermite route gives $\pi A^2\le\tfrac32$, not $\le3$: a factor 2 better than
printed, and \emph{exactly} the bound $(2.5)$ that \S2.2 will reach by a quite
different construction. Either the printed $3$ is a slip for $3/2$, or a different
normalisation of $H_4$ is intended; with the paper's own Fourier convention the
computation gives $3/2$. We record both and change no theorem: $\pi A^2\le3$ is true.
This is also the explicit pair promised in \S1.1.
\end{filled}

\subsection*{2.2 The family $g_a$}

One may also consider the functions
\begin{equation}\label{eq:21}
 g_a(x)=a\gamma(ax)+\gamma\Bigl(\frac xa\Bigr)-(1+a)\gamma(x),\qquad a>1,
\end{equation}
which satisfy the conditions of \eqref{eq:11}. Then any expression of the form
\begin{equation}\label{eq:22}
 \int_1^\infty g_a(x)\,d\tau(a)
\end{equation}
where $\tau$ is a measure on $]1,\infty[$ such that the integral converges for every
$x$ and is $\ge0$ at infinity, is a candidate function. (It seems difficult to
determine a simple characteristic property of $\tau$ ensuring that \eqref{eq:22}
converges and is positive at infinity.)

\begin{expansion}
\textbf{$g_a$ is admissible.} \Ledger{C4} With the paper's convention,
$\widehat{\gamma(a\,\cdot)}(y)=a^{-1}\gamma(y/a)$. Hence
\[\widehat{g_a}=a\cdot a^{-1}\gamma(y/a)+a\gamma(ay)-(1+a)\gamma(y)=g_a(y),\]
so $g_a$ is self-dual; and $g_a(0)=a+1-(1+a)=0$. It is real, even and integrable, so
it is a competitor in \eqref{eq:11} provided it is not identically zero, which
$g_a''(0)\ne0$ shows for $a\ne1$.
\end{expansion}

We first study $A(g_a)$. It is convenient to set $X=\pi x^2$ and $G_a(X)=g_a(x)$, so
\[G_a(X)=a\,e^{-a^2X}+e^{-a^{-2}X}-(1+a)e^{-X} .\]
Moreover
\begin{equation}\label{eq:23}
 H_a(X)=e^XG_a(X)=a\,e^{(1-a^2)X}+e^{(1-a^{-2})X}-1-a
\end{equation}
is a convex function satisfying
\[H_a(0)=0,\qquad H_a'(0)=-a^{-2}(a^2-1)(a^3-1)<0\]
and tending to $+\infty$ with $X$. It therefore admits a unique zero $X_a>0$, and
\[A(g_a)=\sqrt{\frac{X_a}{\pi}} .\]

\begin{expansion}
\textbf{Convexity, the derivative, and uniqueness.} \Ledger{C6--C10}
$H_a$ is a sum of two exponentials $e^{cX}$ (with $c=1-a^2<0$ and $c=1-a^{-2}>0$) and
a constant, hence convex. $H_a(0)=a+1-1-a=0$. Differentiating,
$H_a'(0)=a(1-a^2)+(1-a^{-2})$; and
\[-a^{-2}(a^2-1)(a^3-1)=-a^{-2}(a^5-a^3-a^2+1)=-a^3+a+1-a^{-2},\]
which is the same. For $a>1$ both factors $(a^2-1),(a^3-1)$ are positive, so
$H_a'(0)<0$. A convex function with $H(0)=0$, $H'(0)<0$ and $H(+\infty)=+\infty$ is
negative on $(0,X_a)$ and positive on $(X_a,\infty)$ for exactly one $X_a>0$: it
decreases then increases, so it has exactly two zeros counted with the one at the
origin. Since $X=\pi x^2$ is increasing in $|x|$ and $e^X>0$, the last sign change of
$g_a$ is at $|x|=\sqrt{X_a/\pi}$.
\end{expansion}

It is natural to study the variation of $X_a$, first for $a$ near $1$. Putting
$a=1+h$, $h>0$, one finds for fixed $X$, modulo $O(h^5)$, an expansion
$P_1h+P_2h^2+P_3h^3+P_4h^4$, where
\[P_1=0,\quad P_2=2X(2X-3),\quad P_3=-X(2X-3),\quad
  P_4=-5X+15X^2-\tfrac{28}{3}X^3+\tfrac43X^4 .\]

\begin{ournotation}
\Ledger{C11} The original writes the exponent of the second term as
$e^{X(2h-3h^2+3h^3-4h^4)X}$, with a doubled $X$, and lists the series as
$P_1h+P_2h^2+P_3h^3+P_4h^5+O(h^5)$, where $P_4h^4$ is meant. Both are typesetting
slips; the four polynomials as printed are correct, and we have recomputed all of
them. \Ledger{C12}
\end{ournotation}

The expression for $P_2$ shows that for $h$ small enough $H_a(X)>0$ if $X>\tfrac32$
and $H_a(X)<0$ if $X<\tfrac32$. Consequently
\begin{equation}\label{eq:24}
 \lim_{a\to1^+}X_a=\tfrac32,
\end{equation}
which furnishes an explicit bound
\begin{equation}\label{eq:25}
 A\le\sqrt{\frac{3}{2\pi}} .
\end{equation}

But this simple bound cannot be the true value of $A$. For $X=\tfrac32$, $P_2$ and
$P_3$ vanish, and $P_4(\tfrac32)=\tfrac32$. So for $h$ small and non-zero,
$X_a<\tfrac32$.

\begin{expansion}
\Ledger{C13, C14} Both checked: $\sqrt{3/(2\pi)}=0{,}6909883$, and
$P_4(3/2)=-7{,}5+33{,}75-31{,}5+6{,}75=1{,}5$. The inference is that at
$X=\tfrac32$ the expansion begins at order $h^4$ with a \emph{positive} coefficient,
so $H_a(\tfrac32)>0$ for small $h\ne0$; since $H_a$ is negative exactly on
$(0,X_a)$, this forces $X_a<\tfrac32$. Note this makes \eqref{eq:25} a limit that is
approached but not attained, and identifies the direction of improvement.
\end{expansion}

If $a\to+\infty$ then $X_a\to+\infty$; in fact $X_a=\log a+O(1)$.

\begin{filled}
\Ledger{C15} The paper says ``un calcul simple montre''. Here it is. In
\eqref{eq:23} with $a$ large, $ae^{(1-a^2)X}\to0$ for $X$ bounded away from $0$, so
the equation $H_a(X)=0$ reads $e^{(1-a^{-2})X}=1+a+o(1)$, i.e.
$X(1-a^{-2})=\log(1+a)+o(1)$, i.e. $X_a=\log a+O(1)$ since $1-a^{-2}\to1$ and
$\log(1+a)=\log a+O(1)$.
\end{filled}

We have not determined the minimal value of $X_a$, but it is easy to estimate it
semi-heuristically. The value $a=\sqrt2$ gives, on setting $q=e^{\frac12X_a}$,
\[q^3-(1+\sqrt2)q^2+\sqrt2=0;\]
if $q\ne1$, the quadratic equation $q^2-\sqrt2q-\sqrt2=0$ gives
$q=\tfrac{\sqrt2}2(1+\sqrt{1+2\sqrt2})$ and
\[X_a=2\log q=1{,}4749\ldots\qquad(a=\sqrt2).\]
The value $a=2$ gives, for $q=e^{\frac34X}$,
\[q^4-2\,\frac{q^4-1}{q-1}=0 .\]
The solution $q>1$ is $q=2{,}9744\ldots$, whence
\[X_a=1{,}4534\ldots\qquad(a=2).\]

\begin{expansion}
\textbf{Both reductions, derived; both values, recomputed.} \Ledger{C16, C17}

\emph{$a=\sqrt2$.} Then $a^2=2$, $a^{-2}=\tfrac12$, and \eqref{eq:23} is
$H(X)=\sqrt2e^{-X}+e^{X/2}-1-\sqrt2$. With $q=e^{X/2}$, $e^{-X}=q^{-2}$, so
$H=0$ reads $\sqrt2q^{-2}+q-1-\sqrt2=0$; multiplying by $q^2$,
\[q^3-(1+\sqrt2)q^2+\sqrt2=0 .\]
$q=1$ is a root ($1-1-\sqrt2+\sqrt2=0$), and dividing out $(q-1)$ leaves
$q^2-\sqrt2q-\sqrt2=0$ --- the step the paper omits. Its positive root is
$q=\tfrac{\sqrt2}{2}\bigl(1+\sqrt{1+2\sqrt2}\bigr)$, and $X=2\log q$.
\emph{Recomputed: $q=2{,}090657851$, $X_a=1{,}474957555$.}

\emph{$a=2$.} Then $a^2=4$, $a^{-2}=\tfrac14$, and $H(X)=2e^{-3X}+e^{3X/4}-3$. With
$q=e^{3X/4}$, $e^{-3X}=q^{-4}$, so $H=0$ reads $2q^{-4}+q-3=0$, i.e.
$q^5-3q^4+2=0$. Factoring out $(q-1)$: $q^4-2q^3-2q^2-2q-2=0$. The printed form
$q^4-2(q^4-1)/(q-1)=0$ is the same equation, since
$(q^4-1)/(q-1)=q^3+q^2+q+1$. \emph{Recomputed: $q=2{,}974449245$,
$X_a=\tfrac43\log q=1{,}453411859$.}
\end{expansion}

It is likely that this is roughly the optimal value accessible by this method.
Indeed if one solves $H_a(X)=0$, with $H_a$ given by \eqref{eq:23}, and supposes
$a\ge2$, the first term is negligible. So $X_a$ is roughly
\[\frac{\log(1+a)}{1-a^{-2}} .\]
The extremum of this expression is attained for $a(a-1)=2\log(1+a)$, which gives
\[a=2{,}08137\ldots\]

\begin{expansion}
\textbf{The printed condition, and what the number is.} \Ledger{C18, C19}
The original prints $a(1-a)=2\log(1+a)$. For $a>1$ the left side is negative and the
right side positive, so that equation has no root there; the intended condition is
$a(a-1)=2\log(1+a)$, whose root is $a=2{,}081376974$ --- exactly the printed value.
So the equation is mistyped and the number is right.

\emph{But it is the extremum of the surrogate, not of $X_a$.} Recomputing the true
minimiser of $X_a$ over $a>1$ on a fine grid gives $a\approx1{,}89$ with
$X\approx1{,}45200$, against $X_{2}=1{,}45341$ and $X_{2{,}08137}=1{,}45623$. The
value $2{,}08137$ minimises the large-$a$ approximation
$\log(1+a)/(1-a^{-2})$, not $X_a$ itself. The paper's hedging (``vraisemblable'',
``\`a peu pr\`es'') is fair; we draw the distinction because the two numbers are
easily confused.
\end{expansion}

\subsection*{2.3 The correction, and Theorem 2}

In all cases the minimal value $A(g_a)$ so obtained is not the value \eqref{eq:11}
sought. Consider indeed $a_0$ such that $X_0=X_{a_0}$ is minimal, and $H_0=H_{a_0}$,
positive on $[X_0,\infty)$.

Let $a$ (for instance near 1) be such that $X_a>X_0$. On $[X_a,\infty)$, $H_a$ is
$\ge0$ and its order of growth as $X\to+\infty$, namely $e^{(1-a^{-2})X}$, is smaller
than that of $H_{a_0}$ if $a<a_0$. There is therefore $T>0$ such that
$H_{a_0}-TH_a$ is $\ge0$ on $[X_a,\infty)$. But this function is $>0$ on
$[X_0,X_a)$, hence on a neighbourhood of $X_0$, hence for $X\ge X'$ with $X'<X_0$.

The same argument applies taking any $a_0$ such that $X_0<\tfrac32$. For $a_0=2$ one
can determine the optimal correction (corresponding to $a$ very close to 1), giving a
function $\ge0$ on $[X'',\infty)$, with
\begin{equation}\label{eq:26}
 X''=1{,}25\ldots,\qquad A\le0{,}63\ldots
\end{equation}
We have made only a very rough calculation. We state the result nonetheless, to be
compared with Theorem 1.

\begin{thmB}
$A\le0{,}64$ and $B_1\le0{,}41$.
\end{thmB}

\begin{expansion}
\textbf{The correction argument, unpacked.} \Ledger{C20} Three claims are compressed
into that paragraph.

\emph{(a) $H_{a_0}-TH_a$ is a legitimate competitor.} $g_{a_0}-Tg_a$ is a real even
$L^1$ function, self-dual and vanishing at $0$, being a linear combination of two such
(\S2.2). Its last sign change is at $X$ where $H_{a_0}-TH_a$ last changes sign, since
$G=e^{-X}H$ and $e^{-X}>0$. So any $X'$ with $H_{a_0}-TH_a\ge0$ on $[X',\infty)$
gives $A\le\sqrt{X'/\pi}$.

\emph{(b) Such a $T$ exists.} On $[X_a,\infty)$ both $H_{a_0}$ and $H_a$ are $\ge0$,
and $H_a>0$ on $(X_a,\infty)$. As $X\to\infty$, $H_{a_0}\sim e^{(1-a_0^{-2})X}$ and
$H_a\sim e^{(1-a^{-2})X}$; since $a<a_0$ we have $1-a^{-2}<1-a_0^{-2}$, so
$H_{a_0}/H_a\to+\infty$. The ratio is therefore bounded below by a positive number on
$[X_a,\infty)$ --- it is continuous and positive on the interior and tends to
$+\infty$ at both ends of the relevant range --- and $T:=\inf_{X\ge
X_a}H_{a_0}(X)/H_a(X)>0$ is the largest admissible value.

\emph{(c) The zero moves left.} On $[X_0,X_a)$ we have $H_{a_0}\ge0$ while $H_a<0$
(as $X<X_a$), so $H_{a_0}-TH_a>0$ there. By continuity it stays $>0$ slightly to the
left of $X_0$, so the last sign change $X''$ satisfies $X''<X_0$. That is the gain.
\end{expansion}

\begin{expansion}
\textbf{The numbers in \eqref{eq:26}, recomputed.} \Ledger{C21} Taking $a_0=2$ and
letting $a\to1^+$ --- the ``optimal correction'' the paper describes --- the
construction converges to
\[X''=1{,}1512203,\qquad A\le\sqrt{X''/\pi}=0{,}6053468 .\]
Optimising over $a_0$ as well, $a_0\approx1{,}5$--$1{,}7$ gives $X''\approx1{,}1210$
and $A\le0{,}5973$.

The printed $X''=1{,}25\ldots$, $A\le0{,}63\ldots$ are therefore \emph{weaker} than
what the method yields --- which is exactly what the authors' own caveat says
(``nous n'avons fait qu'un calcul tr\`es approch\'e''). Theorem 2 as stated is true;
we do not restate it. See \S5(2).
\end{expansion}

% =====================================================================
\section*{\S3. Higher dimensions}
% =====================================================================

We work in Euclidean $\R^d$, with $x\cdot y=\sum_1^dx_iy_i$, $\|x\|=(x\cdot x)^{1/2}$,
and
\begin{equation}\label{eq:31}
 \fh(y)=\int f(x)e^{-2i\pi x\cdot y}dx,\qquad
 f(x)=\int\fh(y)e^{2i\pi x\cdot y}dy .
\end{equation}
More generally, on a Euclidean space $E$ of dimension $d$ with invariant measure
normalised so that the cube spanned by an orthonormal basis has measure 1, and $x\cdot
y$ the inner product, \eqref{eq:31} inverts as stated.

\begin{ournotation}
\Ledger{D1} That normalisation of Haar measure on an \emph{abstract} Euclidean space
is stated once here and used again in \S4, where $F_\infty$ carries no preferred
coordinates. We flag it because it is the hinge between \S3 and \S4.
\end{ournotation}

Consider Fourier pairs satisfying
\begin{enumerate}[label=\textbf{\arabic*)},leftmargin=2.4em]
\item $f,\fh$ real and even, not identically zero;
\item $f(0)\le0$ and $\fh(0)\le0$;
\item $f(x)\ge0$ for $\|x\|\ge a_f$, $\fh(y)\ge0$ for $\|y\|\ge a_{\fh}$,
\end{enumerate}
put $A(f)=\inf\{r>0:f(x)\ge0\text{ if }\|x\|>r\}$ and $B_d=\inf A(f)A(\fh)$.

Let $f^\#(x)$ be the (invariant) average of $f$ over the sphere of radius $\|x\|$:
then $\widehat{f^\#}=(\fh)^\#$, and $f^\#$ and $\widehat{f^\#}$ are not zero ---
otherwise $f$ and $\fh$ would have compact support by 3). Since $A(f^\#)\le A(f)$ and
$A(\widehat{f^\#})\le A(\fh)$, one may restrict to radial pairs. Since
$(f(x/\lambda))^{\wedge}=\lambda^d\fh(\lambda y)$, the argument of \S1 applies and
\begin{equation}\label{eq:34a}
 B_d=A^2,\qquad A=\inf A(f),
\end{equation}
\emph{the infimum over $f\in L^1(\R^d)$ radial, not identically zero, with $f=\fh$ and
$f(0)=0$}, having added if necessary a multiple of the radial self-dual
$\gamma(x)=e^{-\pi\|x\|^2}$.

\begin{filled}
\textbf{The radialisation, both halves.} \Ledger{D3} The paper compresses this to one
sentence.

\emph{Why $f^\#$ is a competitor.} Spherical averaging commutes with the Fourier
transform on radial-invariant data: $\widehat{f^\#}=(\fh)^\#$, because the transform
commutes with rotations and averaging is an integral over the rotation group. If $f\ge
0$ for $\|x\|\ge a_f$ then so is its spherical average, giving $A(f^\#)\le A(f)$; and
$f^\#(0)=f(0)\le0$.

\emph{Why $f^\#\not\equiv0$.} Suppose $f^\#\equiv0$. By 3), $f\ge0$ on
$\{\|x\|\ge a_f\}$, and a non-negative function with vanishing spherical averages
vanishes a.e.\ there. So $f=0$ a.e.\ on $\{\|x\|\ge a_f\}$: $f$ has compact support.
The same argument applied to $\fh$ gives $\fh$ compactly supported. By Background
\ref{bg:cpt}, $f\equiv0$ --- contradicting 1). The paper states the first half and
leaves the contradiction implicit.
\end{filled}

\begin{expansion}
\textbf{What changes in \eqref{eq:34a} relative to \S1.2.} \Ledger{D4} Only the
dilation exponent: $(f(x/\lambda))^{\wedge}=\lambda^d\fh(\lambda y)$, so
$A(f)\mapsto\lambda A(f)$ and $A(\fh)\mapsto\lambda^{-1}A(\fh)$ exactly as before ---
the $\lambda^d$ multiplies the \emph{values}, not the argument, and $A$ is unchanged
by positive scalar multiples. Hence the product is again dilation-invariant, the
normalisation $A(f)=A(\fh)$ is again available, and the rest of \S1.2 --- the
admissibility of $f+\fh$, the compact-support contradiction, the Gaussian correction
--- is verbatim with $|x|$ replaced by $\|x\|$ and $\gamma(x)=e^{-\pi\|x\|^2}$.
\end{expansion}

\begin{thmC}
$\displaystyle B_d\ge\frac1\pi\Bigl(\tfrac12\Gamma\Bigl(\tfrac d2+1\Bigr)\Bigr)^{2/d}
 >\frac{d}{2\pi e}$.
\end{thmC}

\begin{proof}
It is modelled on the case $d=1$, replacing the interval $(-A(f),A(f))$ by the ball of
centre $O$ and radius $A(f)$, whose volume (which is $\ge\tfrac12$) is
$\frac1{\Gamma(\frac d2+1)}(A(f))^d\pi^{d/2}$.
\end{proof}

\begin{filled}
\textbf{The proof, written out, and the second inequality, proved.}
\Ledger{D7, D8}

\emph{First inequality.} Normalise $\int_{\R^d}|f|=1$. As in Theorem 1:
$\int f=\fh(0)=0$ gives $\int f^+=\int f^-=\tfrac12$; $f\ge0$ outside the ball
$B(0,A)$ gives $\int_{\|x\|\ge A}|f|=\int_{\|x\|\ge A}f^+\le\tfrac12$, hence
$\int_{B(0,A)}|f|\ge\tfrac12$; and $\|f\|_\infty\le\|\fh\|_1=1$. Therefore
$\operatorname{vol}B(0,A)\ge\tfrac12$. With
$\operatorname{vol}B(0,A)=\pi^{d/2}A^d/\Gamma(\frac d2+1)$ this reads
$A^d\ge\Gamma(\frac d2+1)/(2\pi^{d/2})$, i.e.
\[B_d=A^2\ge\frac1\pi\Bigl(\tfrac12\Gamma\bigl(\tfrac d2+1\bigr)\Bigr)^{2/d}.\]

\emph{Second inequality.} It is equivalent to
$\bigl(\tfrac12\Gamma(m+1)\bigr)^{1/m}>m/e$ with $m=d/2$, i.e.\ to
$\Gamma(m+1)>2(m/e)^m$. By Background \ref{bg:stir},
$\Gamma(m+1)>\sqrt{2\pi m}\,(m/e)^m$, so it suffices that $\sqrt{2\pi m}\ge2$, i.e.
$m\ge2/\pi=0{,}6366$, i.e. $d\ge4/\pi=1{,}2732$ --- so for every $d\ge2$. For $d=1$
check directly: $\Gamma(3/2)=\sqrt\pi/2=0{,}8862269$ and
$2(1/2e)^{1/2}=0{,}8577639$, so the inequality holds there too. The paper asserts the
inequality without argument. \emph{Verified numerically for
$d=1,2,3,4,8,12,24,48,100$ as well.}
\end{filled}

Setting $X=\pi\|x\|^2$, the argument of \S2 leads to the natural functions
$g_a(x)=G_a(X)$ with
\[G_a(X)=a^de^{-Xa^2}+e^{-Xa^{-2}}-(1+a^d)e^{-X},\qquad
  H_a(X)=a^de^{(1-a^2)X}+e^{(1-a^{-2})X}-(1+a^d),\ a>1 .\]
\begin{expansion}
\textbf{$g_a$ is still self-dual in dimension $d$.} \Ledger{D9} With
$\gamma(x)=e^{-\pi\|x\|^2}$ one has $\widehat{\gamma(a\,\cdot)}(y)=a^{-d}\gamma(y/a)$,
so
\[\widehat{a^d\gamma(a\,\cdot)}=\gamma(\cdot/a),\qquad
  \widehat{\gamma(\cdot/a)}=a^d\gamma(a\,\cdot),\]
and the two outer terms of $g_a$ swap while $(1+a^d)\gamma$ is fixed. Hence
$\widehat{g_a}=g_a$, and $g_a(0)=a^d+1-(1+a^d)=0$. This is why the coefficient is
$a^d$ here where it was $a$ in \eqref{eq:21}: it is exactly the Jacobian factor.
\end{expansion}

It is convenient to set $a^2=1+k$, $d=2c$, whence
\[H_a(X)=(1+k)^ce^{-kX}+e^{(1-(1+k)^{-1})X}-1-(1+k)^c .\]
The derivative at the origin in $X$ is
$\frac{k}{1+k}\bigl(1-(1+k)^{c+1}\bigr)<0$; the convexity argument of \S2 shows that
$H_a$ has a unique positive zero $X_a$. As before we compute an expansion in $k$ to
order 4:
$H_a(X)=P_1k+P_2k^2+P_3k^3+P_4k^4+O(k^5)$ with
\[P_1=0,\quad P_2=X(X-c-1),\quad P_3=\tfrac12(c-2)X(X-c-1),\]
\[P_4=\tfrac1{12}X\bigl\{X^3-(2c+6)X^2+(3c(c-1)+18)X-(2c(c-1)(c-2)+12)\bigr\}.\]

\begin{expansion}
\Ledger{D10, D11} The derivative at the origin: differentiating the displayed form at
$X=0$ gives $-k(1+k)^c+(1-(1+k)^{-1})=-k(1+k)^c+\frac{k}{1+k}$, which is
$\frac{k}{1+k}\bigl(1-(1+k)^{c+1}\bigr)$; for $k>0$ and $c>0$ the bracket is negative.
All four polynomials \emph{recomputed and confirmed} at several $(X,c)$.
\end{expansion}

As in dimension 1, $P_2$ and $P_3$ vanish for
\begin{equation}\label{eq:Xd}
 X=X(d):=\tfrac d2+1 .
\end{equation}
Moreover $P_2>0$ for $X>X(d)$ and $<0$ for $X<X(d)$. Letting $k\to0$ we deduce
$\lim_{a\to1}X_a=\tfrac d2+1$.

\begin{ournotation}
\Ledger{D5} The original labels this equation $(3.4)$, which it has already used for
$B_d=A^2$. We have relabelled; the second occurrence is \eqref{eq:Xd} here.
\end{ournotation}

To understand the position of $X_a$ relative to $X(d)$ as $a\to1$, compute
$Q_4(X(d))$, where $P_4=\tfrac X{12}Q_4$. The computation gives
\[Q_4(c+1)=-c^2+1 .\]
For $d>2$ this term is $<0$, so $H_a(X(d))<0$ for $a$ close to 1, which shows
\[X_a>\tfrac d2+1\qquad(a>1\text{ close enough to }1).\]
It is therefore possible that the value \eqref{eq:Xd} is optimal. For $d=1$ it is
not, as \S2 showed.

\begin{expansion}
\Ledger{D12, D13} $Q_4(c+1)=-c^2+1$ \emph{recomputed and confirmed} for
$c=\tfrac12,1,\tfrac32,3,5$. The inference: at $X=X(d)$ the $k$-expansion starts at
order $k^4$ with coefficient $\tfrac{X}{12}(-c^2+1)$, negative for $c>1$ i.e. $d>2$.
So $H_a(X(d))<0$ for small $k>0$, and since $H_a<0$ exactly on $(0,X_a)$, this forces
$X_a>X(d)$. Note this is the \emph{opposite} direction to $d=1$ (where
$P_4(3/2)>0$ forced $X_a<3/2$): in high dimension the family $g_a$ does no better
than its limit, which is why the authors suspect $X(d)$ is optimal.
\end{expansion}

For $d=2$, $Q_4(c+1)=0$, so we must compute the expansion of
\begin{equation}\label{eq:35}
 H_a(2)=(1+k)e^{-2k}+e^{2(1-\frac1{1+k})}-2-k
\end{equation}
to order at least 5. The Taylor expansion at $0$ of
$f(z)=e^{2(1-\frac1{1+z})}=e^{2\frac{z}{1+z}}=\sum_0^\infty q_nz^n$ is computed by the
residue theorem: putting $w=\frac z{1+z}$, $z=\frac w{1-w}$,
$dz=\frac{dw}{(w-1)^2}$, one has
\[q_n=\Res_{z=0}\frac{f(z)}{z^{n+1}}
 =\frac1{2i\pi}\oint e^{\frac{2z}{1+z}}\frac{dz}{z^{n+1}}
 =\frac1{2i\pi}\oint e^{2w}\frac{(1-w)^{n+1}}{w^{n+1}}\frac{dw}{(1-w)^2}
 =\Res_{w=0}\frac{(1-w)^{n-1}}{w^{n+1}}e^{2w} .\]
\begin{expansion}
\textbf{The substitution in the residue.} \Ledger{D15} With $w=z/(1+z)$ the map is a
biholomorphism near $0$ fixing $0$, so residues transport; $z=w/(1-w)$ gives
$z^{-(n+1)}=(1-w)^{n+1}w^{-(n+1)}$ and $dz=(1-w)^{-2}dw$, whence the exponent
$n+1-2=n-1$ in the last display. The point of the change of variable is that
$e^{2z/(1+z)}$ becomes $e^{2w}$, whose Taylor coefficients are immediate.
\end{expansion}

In particular $p_5$ is the sum of
\begin{equation}\label{eq:36}
 \frac{2^4}{4!}-\frac{2^5}{5!}
\end{equation}
coming from the first term of \eqref{eq:35}, and of the term in $w^5$ of
$e^{2w}(1-w)^4$, equal to
\begin{equation}\label{eq:37}
 \frac{2^5}{5!}-4\cdot\frac{2^4}{4!}+6\cdot\frac{2^3}{3!}-4\cdot\frac{2^2}{2!}+2 .
\end{equation}
One finds $p_5=0$. Likewise $p_6$ is the sum of
\begin{equation}\label{eq:38}
 -\frac{2^5}{5!}+\frac{2^6}{6!}
\end{equation}
and of
\begin{equation}\label{eq:39}
 \frac{2^6}{6!}-5\cdot\frac{2^5}{5!}+10\cdot\frac{2^4}{4!}-10\cdot\frac{2^3}{3!}
 +5\cdot\frac{2^2}{2!}-2,
\end{equation}
whence $p_6=-\tfrac4{45}<0$. For $a$ very close to 1 we therefore have $H_a(2)<0$ and
$X_a>X(2)=2$. Here again it is possible that the bound \eqref{eq:Xd} is optimal.

\begin{ournotation}
\Ledger{D16} The coefficients of this last expansion are called $p_5,p_6$ while those
of the earlier one were $P_1,\dots,P_4$. They are different objects --- these are
Taylor coefficients of $H_a(2)$ in $k$, those were polynomials in $X$ --- and the
near-identical names are the original's.
\end{ournotation}

\begin{expansion}
\Ledger{D17} \emph{Both recomputed exactly.} \eqref{eq:36}$=0{,}4$ and
\eqref{eq:37}$=-0{,}4$, so $p_5=0$. \eqref{eq:38}$=-0{,}17\overline{7}$ and
\eqref{eq:39}$=+0{,}0\overline{8}$, so $p_6=-0{,}0\overline{8}=-\tfrac4{45}$. A
direct Taylor expansion of \eqref{eq:35} confirms $p_0=\dots=p_5=0$ and
$p_6=-4/45$.
\end{expansion}

To conclude this section, note that for every $d\ge2$ we have obtained the upper bound
\begin{equation}\label{eq:310}
 B_d\le\Bcal_d\le\frac{d+2}{2\pi},
\end{equation}
where $\Bcal_d$ is defined, as in \S1, by functions of $\Ss(\R^d)$. Moreover the
considerations at the end of \S1, on the bounds for $L^1$ and for $\Ss$, apply. In the
proof of \eqref{eq:16} one must take $T=\delta_b+\delta_{-b}+2\delta_0$ with
$\|b\|<a=A(f)$ and $f(b)<0$; $\widehat T=2\cos(2\pi b\cdot y)+2$ is a positive plane
wave. The rest of the argument is identical, replacing $h+\widehat h$ by the spherical
average of $h+\widehat h$ if one wishes to stay with radial functions. In conclusion,
to be compared with Theorem 3:

\begin{thmD}
\[B_d\le\Bcal_d\le\frac{d+2}{2\pi},\qquad B_d\ge\tfrac12\Bcal_d .\]
\end{thmD}

\begin{expansion}
\Ledger{D18, D19} \eqref{eq:310} is $A^2\le X(d)/\pi=(\tfrac d2+1)/\pi=(d+2)/2\pi$,
using $\lim_{a\to1}X_a=X(d)$ and $A\le A(g_a)$ for every $a$.

The $d$-dimensional version of the $T$ argument: $b\in\R^d$ is now a vector with
$\|b\|<a$, and $\widehat T(y)=2\cos(2\pi b\cdot y)+2\ge0$ is constant on hyperplanes
orthogonal to $b$ --- a plane wave, hence not radial, which is why the spherical
average is taken at the end. Every other step of \S1.4 goes through verbatim with
$|x|$ replaced by $\|x\|$: $g=T*f\ge0$ on $\{\|x\|\ge2a\}$ because each of
$x-b,x,x+b$ then has norm $\ge a$; $\widehat g=\widehat T\fh\ge0$ on
$\{\|y\|\ge a\}$; and the dilation $\lambda=1/\sqrt2$ balances $2a$ against $a$ at
$a\sqrt2$. The conclusion $B_d\ge\tfrac12\Bcal_d$ is \eqref{eq:16} rearranged.
\end{expansion}

% =====================================================================
\section*{\S4. An arithmetic argument}
% =====================================================================

\begin{caution}
\textbf{This section imports heavily, and we have verified none of the imports.}
\Ledger{E1, E5, E8, E16, E21, E23} It rests on Tate's thesis for the functional
equation \eqref{eq:43}, on Armitage for a number field with $\zeta_F(1/2)=0$, on Serre
for the construction of that field, on Odlyzko for discriminant lower bounds, and on
Golod--Shafarevi\v c/Brumer (via Roquette) for infinite class field towers. We state
each precisely and check how the paper uses it, but we have not read those sources.
Everything in \S4 is therefore conditional on them in this digestion.
\end{caution}

\subsection*{4.1 Setting}

Let $F$ be a number field of degree $d$ over $\Q$; $v$ its places, $F_v$ the
completions, $\mathcal O_v$ the ring of integers and $q_v$ the residue cardinality at
a finite place; $\mathbb A_F=\prod'_vF_v$ the adeles, $I_F=\mathbb A_F^\times$ the
ideles, $|x|=\prod_v|x|_v$ the idele norm, and
\[I_F^1=\{x:|x|=1\},\qquad I_F^+=\{x:|x|\ge1\} .\]
The measure $dx=\prod dx_v$ is Haar; at finite places $dx_v$ is Tate's self-dual
measure, at a real place Lebesgue, at a complex place $dz=2\,dx\,dy$. At a complex
place with parameters $z=x+iy$, $w=\xi+i\eta$, Tate defines
$\fh(w)=\int f(z)e^{-2i\pi\Tr(zw)}dz$ with $\Tr(zw)=2\operatorname{Re}(zw)$.

\begin{ournotation}
\textbf{Reconciling the two conventions.} \Ledger{E2} The paper says in one sentence
that for radial functions Tate's complex-place transform ``revient \`a'' the
transform of \S3 with the inner product $z\cdot w=2(x\xi+y\eta)$. Spelled out:
$\Tr(zw)=2(x\xi-y\eta)$ carries a minus sign, but for functions even in each variable
the sign is immaterial, so one may use $z\cdot w=2(x\xi+y\eta)$, which is exactly the
inner product making $dz=2dxdy$ self-dual in the sense of \S3's normalisation. This
is the point at which \S3's abstract-Euclidean convention is needed.
\end{ournotation}

Let $f$ be the Schwartz function on $\mathbb A_F$
\begin{equation}\label{eq:41}
 f(x)=\prod_{v\mid\infty}f_v(x_v)\prod_{v\text{ finite}}f_v^0(x_v)
\end{equation}
with $f_v^0=\mathbf 1_{\mathcal O_v}$. The associated Tate zeta function is
$Z(f,s)=\int_{I_F}f(x)|x|^s\,d^\times x$ for $\operatorname{Re}(s)>1$. Rather than the
decomposable functions of \eqref{eq:41} we consider functions
\begin{equation}\label{eq:42}
 f(x)=f_\infty(x_\infty)\prod_{v\text{ finite}}f_v^0(x_v),\qquad
 f_\infty\in\Ss(\R^d),
\end{equation}
where $\R^d=F_\infty$ is Euclidean via
$\|x_\infty\|^2=\sum_{v\text{ real}}|x_v|^2+\sum_{v\text{ complex}}2\|z_v\|^2$.
Tate's conditions $(z_1),(z_2),(z_3)$ hold for such functions. By Tate,
\begin{equation}\label{eq:43}
 Z(f,s)=\int_{I_F^+}f(x)|x|^sd^\times x+\int_{I_F^+}\fh(x)|x|^{1-s}d^\times x
       +\kappa\frac{\fh(0)}{s-1}-\kappa\frac{f(0)}{s},
\end{equation}
where $\kappa=\dfrac{2^{r_1}(2\pi)^{r_2}hR}{\sqrt{D_F}\,w}$ is the residue of
$\zeta_F$ at $s=1$, $D_F$ the absolute discriminant, $d=r_1+2r_2$. Both integrals in
\eqref{eq:43} converge absolutely for all $s\in\C$.

\subsection*{4.2 The two lemmas}

\begin{lemA}
Let $s$ be a zero of $\zeta_F(s)$ with $\operatorname{Re}(s)>0$. Then $Z(f,s)$
vanishes at $s$, for every choice of $f_\infty\in\Ss(F_\infty)$.
\end{lemA}

Indeed one may write, first for $\operatorname{Re}s>1$,
$Z(f,s)=Z(f_\infty,s)\zeta_F(s)$. Since $Z(f,s)$, $\zeta_F(s)$ and $Z(f_\infty,s)$ are
holomorphic for $s\ne1$, $\operatorname{Re}(s)>0$, the lemma follows.

\begin{filled}
\textbf{The factorisation and the continuation.} \Ledger{E7}

\emph{The factorisation.} For $\operatorname{Re}s>1$ the idele integral factors over
places because $f$ does: $Z(f,s)=\prod_vZ(f_v,s)$. The finite places contribute
$\prod_{v\text{ finite}}(1-q_v^{-s})^{-1}=\zeta_F(s)$, since
$f_v^0=\mathbf 1_{\mathcal O_v}$ and the local zeta integral of the characteristic
function of the integers is the local Euler factor (with $d^\times x_v$ normalised by
the factor $(1-q_v^{-1})^{-1}$ that the paper records). The archimedean places
contribute $Z(f_\infty,s)$ by definition. Hence
$Z(f,s)=Z(f_\infty,s)\zeta_F(s)$ on $\operatorname{Re}s>1$.

\emph{The continuation.} $Z(f,s)$ extends holomorphically to $\{\operatorname{Re}s>0\}
\setminus\{1\}$ by \eqref{eq:43}, whose two integrals are entire. $Z(f_\infty,s)$ is
holomorphic there because $f_\infty\in\Ss$ makes
$\int_{F_\infty}f_\infty\prod|x_v|_v^{\sigma-1}dx$ absolutely convergent for
$\sigma>0$ --- which is the paper's remark about condition $(z_3)$. Two holomorphic
functions agreeing on $\operatorname{Re}s>1$ agree on the connected set
$\{\operatorname{Re}s>0\}\setminus\{1\}$, so the identity persists, and at a zero $s$
of $\zeta_F$ with $\operatorname{Re}s>0$ we get $Z(f,s)=0$.
\end{filled}

For every finite place, $\widehat{f_v^0}=|\mathfrak d_v|^{-1/2}\mathbf
1_{\mathfrak d_v^{-1}}$, where $\mathfrak d_v$ is the different; and
$\prod_{v\text{ finite}}|\mathfrak d_v|=|D_F|$. Consider the first integral of
\eqref{eq:43}. If $f(x)\ne0$ at $x=(x_\infty,x_f)$ then the description of $f_f$ gives
$|x_f|\le1$; since $|x_\infty x_f|\ge1$ on $I_F^+$,
\begin{equation}\label{eq:45}
 |x_\infty|=\prod_{v\mid\infty}|x_v|\ge1 .
\end{equation}
In the second integral, noting $|x_v|\le|\mathfrak d_v|$ if
$x_v\in\mathfrak d_v^{-1}$, one has $|x_f|\le\prod_v|\mathfrak d_v|=|D_F|$, whence
\begin{equation}\label{eq:46}
 |x_\infty|\ge D_F^{-1} .
\end{equation}

\begin{expansion}
\textbf{The two inferences \eqref{eq:45}--\eqref{eq:46}.} \Ledger{E9} On $I_F^+$ the
idele norm satisfies $|x|=|x_\infty|\,|x_f|\ge1$. In the first integral $f\ne0$ forces
$x_v\in\mathcal O_v$ at every finite place, so $|x_v|_v\le1$ and $|x_f|\le1$;
combined with $|x_\infty||x_f|\ge1$ this gives $|x_\infty|\ge1$. In the second, $\fh\ne0$
forces $x_v\in\mathfrak d_v^{-1}$, where $|x_v|_v\le|\mathfrak d_v|$, so
$|x_f|\le\prod_v|\mathfrak d_v|=|D_F|$ and therefore $|x_\infty|\ge|D_F|^{-1}$. The
asymmetry between the two --- $1$ against $D_F^{-1}$ --- is the entire arithmetic
input of \S4: it is what puts the discriminant into Proposition 1.
\end{expansion}

\begin{lemB}
Suppose there is a Fourier pair $(f,\fh)$ on $F_\infty=\R^d$ such that
$f(x_\infty)\ge0$ if $|x_\infty|\ge1$, $f$ takes strictly positive values in a
neighbourhood of $1$ inside $\{|x_\infty|\ge1\}$, $\fh(y_\infty)\ge0$ if
$|y_\infty|\ge D_F^{-1}$, and $f(0)=\fh(0)=0$. Then $\zeta_F(s)\ne0$ for every $s$ in
the interval $]0,1[$.
\end{lemB}

Indeed \eqref{eq:43} is then reduced to its integral terms; $|x|^s$ is strictly
positive on the domain of integration, and the integral \eqref{eq:45} is strictly
positive by the property imposed on $f$. So $Z(f,s)>0$ and $\zeta_F(s)\ne0$ by
Lemma 1.

\begin{expansion}
\textbf{Which hypothesis does what.} \Ledger{E11}
\begin{itemize}[leftmargin=1.4em,itemsep=2pt]
\item $f(0)=\fh(0)=0$ kills the two polar terms $\kappa\fh(0)/(s-1)$ and
  $-\kappa f(0)/s$ of \eqref{eq:43}. Without this, $Z(f,s)$ would have a pole at
  $s=1$ and no sign could be read off.
\item $f\ge0$ on $\{|x_\infty|\ge1\}$ makes the first integral $\ge0$: by
  \eqref{eq:45} the integrand is supported where $|x_\infty|\ge1$.
\item $\fh\ge0$ on $\{|y_\infty|\ge D_F^{-1}\}$ makes the second integral $\ge0$, by
  \eqref{eq:46}.
\item ``$f$ takes strictly positive values near $|x_\infty|=1$'' upgrades the first
  integral from $\ge0$ to $>0$. This is the only hypothesis that gives strictness,
  and strictness is the whole point: $Z(f,s)>0$ contradicts $Z(f,s)=0$.
\item $s\in\,]0,1[$ real makes $|x|^s>0$; for complex $s$ no sign argument is
  available, which is why the conclusion concerns only real zeros.
\end{itemize}
\end{expansion}

Let $x=(x_v)\in F_\infty$. The Euclidean norm compatible with Tate's Fourier
transform is $\|x\|^2=\sum_{v\text{ real}}|x_v|^2+2\sum_{v\text{ complex}}\|x_v\|^2$.
Since $|x|^2=\prod_{v\text{ real}}|x_v|^2\prod_{v\text{ complex}}\|x_v\|^4$, the
arithmetic--geometric inequality gives
\[|x|^{2/d}\le\tfrac1d\|x\|^2 .\]
Putting $r=\|x\|$, $\rho=\|y\|$, one sees that
\[|x|\ge1\ \Longrightarrow\ r\ge\sqrt d,\qquad
  |y|\ge|D_F|^{-1}\ \Longrightarrow\ \rho\ge|D_F|^{-1/d}\sqrt d .\]

\begin{expansion}
\textbf{The AM--GM step, with the places counted.} \Ledger{E12, E13} There are $r_1$
real and $r_2$ complex places, $d=r_1+2r_2$. Write the $d$ non-negative quantities
\[\underbrace{|x_v|^2}_{v\text{ real},\ r_1\text{ terms}},\qquad
  \underbrace{\|x_v\|^2,\ \|x_v\|^2}_{v\text{ complex},\ 2r_2\text{ terms}}\]
--- each complex place contributing its $\|x_v\|^2$ \emph{twice}. Their sum is
exactly $\|x\|^2$ and their product is exactly $|x|^2$. AM--GM on $d$ terms gives
$(|x|^2)^{1/d}\le\frac1d\|x\|^2$, which is the display. The factor 2 in
$\|x\|^2$ and the fourth power in $|x|^2$ are the same phenomenon counted once each,
which is easy to mis-count.

Then $|x|\ge1$ gives $\|x\|^2\ge d|x|^{2/d}\ge d$, i.e. $r\ge\sqrt d$; and
$|y|\ge|D_F|^{-1}$ gives $\rho^2\ge d|D_F|^{-2/d}$, i.e.
$\rho\ge|D_F|^{-1/d}\sqrt d$.
\end{expansion}

\begin{propA}
Suppose there is a number field $F$ of degree $d$ and discriminant $D$ such that
$\zeta_F(s)$ has a zero in $]0,1[$. Then $\Bcal_d\ge d\,|D|^{-1/d}$.
\end{propA}

Conversely, of course, $\zeta_F$ has no real zero if $d\,|D|^{-1/d}>\Bcal_d$.

\begin{filled}
\textbf{The proof, which the paper calls evident.} \Ledger{E14, E15} Suppose
$d|D|^{-1/d}>\Bcal_d$. By definition of $\Bcal_d$ as an infimum over Schwartz
competitors, and by the radial reduction of \S3, there is a radial Schwartz Fourier
pair $(f,\fh)$ with $f=\fh$ up to the dilation of \S1.2, $f(0)=\fh(0)=0$, and
\[A(f)A(\fh)<d\,|D|^{-1/d} .\]
Dilating (which preserves the product) we may arrange
$A(f)\le\sqrt d$ and $A(\fh)\le|D|^{-1/d}\sqrt d$, i.e. $f\ge0$ for $r\ge\sqrt d$ and
$\fh\ge0$ for $\rho\ge|D|^{-1/d}\sqrt d$. By the two implications just displayed,
these are exactly the hypotheses of Lemma 2 in terms of the idele norm. One may
further assume $f>0$ on $\{\sqrt d\le\|x\|\le\sqrt d+\varepsilon\}$ --- otherwise $f$
vanishes on an open set beyond its last sign change, and one shrinks $A(f)$ ---
which supplies the strictness hypothesis, since $\|\mathbf 1\|=\sqrt d$ for the idele
$\mathbf 1$. Lemma 2 then gives $\zeta_F\ne0$ on $]0,1[$. The proposition is the
contrapositive.
\end{filled}

\subsection*{4.3 Fields with a real zero, and linear growth}

It is difficult to find fields $F$ satisfying the hypothesis of the proposition.
However, the decomposition, for $F$ Galois over $E$, of $\zeta_F$ into Artin
$L$-functions for $E$ allowed Armitage to exhibit such a zero (at $s=\tfrac12$ of
course, in accordance with the Riemann hypothesis). More precisely, Armitage considers
an explicit extension $F$ of $E=\Q(\sqrt{3(1+i)})$, of degree 12 over $E$ and hence of
degree 48 over $\Q$, constructed by Serre, and shows that $\zeta_F(\tfrac12)=0$.
Consequently number theory implied a priori the following weak version of Theorem 3:

\begin{propB}
If $d$ is a multiple of 48, then $\Bcal_d$ is strictly positive.
\end{propB}

For $d=48$ this follows from the existence of $F$. Suppose $d=48c$. There is a
cyclotomic extension $L$ of $\Q$ of degree $c$ linearly disjoint from $F$. Then $LF$
has degree $d$ over $\Q$, and $\zeta_F$ divides $\zeta_{LF}$ since $LF/F$ is abelian
and $\zeta_{LF}$ therefore factors into Dirichlet $L$-functions relative to $F$.

\begin{expansion}
\textbf{Two silent uses.} \Ledger{E18}
\emph{(a) A cyclotomic field of each degree, disjoint from $F$.} The cyclotomic fields
$\Q(\zeta_n)$ have degrees $\varphi(n)$, which do not take every value; but one may
take a subfield of $\Q(\zeta_p)$ for a large prime $p\equiv1\pmod c$, whose Galois
group is cyclic of order $p-1$, and hence has a subfield of degree exactly $c$.
Choosing $p$ unramified in $F$ makes $L$ linearly disjoint from $F$.
\emph{(b) $\zeta_F\mid\zeta_{LF}$.} For $LF/F$ abelian with group $G$,
$\zeta_{LF}=\prod_{\chi\in\widehat G}L(s,\chi)$ over $F$, and the trivial character
contributes $\zeta_F$. So the quotient is a product of $L$-functions, entire, and
every zero of $\zeta_F$ is a zero of $\zeta_{LF}$. Proposition 1 applied to $LF$ then
gives $\Bcal_d>0$.
\end{expansion}

One may ask whether Proposition 1 implies a restriction on the discriminants of fields
whose $\zeta_F$ has a real zero. In that case
\begin{equation}\label{eq:47}
 |D|^{1/d}\ge\frac{d}{\Bcal_d} .
\end{equation}
But unconditionally, by Theorem 3, $\dfrac{d}{\Bcal_d}<2\pi e=17{,}079\cdots$,
whereas Odlyzko's lower bounds give in general $|D|^{1/d}\ge22{,}2\,(1+o(1))$ as
$d\to\infty$. Consequently \eqref{eq:47} is automatically satisfied, at least for $d$
large enough.

\begin{expansion}
\Ledger{E20, E21, E22} The chain is $\Bcal_d\ge B_d\ge d/(2\pi e)$ from
\eqref{eq:14} and Theorem 3, so $d/\Bcal_d\le2\pi e=17{,}0795$. Since Odlyzko's
constant $22{,}2$ exceeds it, \eqref{eq:47} carries no information. The original
prints Odlyzko's bound as $22{,}2(1+0(d))$; $o(1)$ is meant.

This is a \emph{negative} result, and the authors include it deliberately:
Proposition 2 does not improve on Theorem 3. What follows is the part that does say
something.
\end{expansion}

Proposition 2 therefore does not lead to an interesting lower bound for $\Bcal_d$. It
is striking to note nonetheless that, for certain degrees at least, number theory
implied the linear growth in $d$ given by Theorem 3.

Let $p$ be a prime. By the theorems of Golod--Shafarevi\v c and Brumer there is a
sequence of fields $E_p^1\subset E_p^2\subset\cdots$ where $E_p^1$, of degree $p(p-1)$
over $\Q$, is an extension of degree $p$ of $\Q(\zeta_p)$, and $E_p^{n+1}/E_p^n$ is
abelian, unramified, of degree $p$. Consider the extensions $F_i=F\,E_p^i$ of $F$;
extracting a strictly increasing subsequence gives
$F_0\subset F_1\subset\cdots\subset F_m\subset\cdots$, each extension abelian of
degree $p$. In the absence of relative ramification, a classical formula gives
\begin{equation}\label{eq:48}
 D(F_m)=D(F_0)^{p^m}:=D^{p^m} .
\end{equation}
The successive extensions from $F$ being abelian, $\zeta_F$ divides $\zeta_{F_m}$ for
all $m$. Proposition 1 then gives, for $d=d_0p^m$ with $d_0=[F_0:\Q]$,
\begin{equation}\label{eq:49}
 \Bcal_d\ge C\,d,\qquad C=|D|^{-1/d_0} .
\end{equation}
For such sequences of degrees, \eqref{eq:310} and \eqref{eq:49} therefore show that
the growth of $\Bcal_d$ --- and hence of $B_d\ge\tfrac12\Bcal_d$ --- is linear in $d$.

\begin{expansion}
\textbf{Why \eqref{eq:49} follows, and two printed slips.} \Ledger{E24, E25, E26}
By Proposition 1 applied to $F_m$, $\Bcal_d\ge d\,|D(F_m)|^{-1/d}$ with $d=d_0p^m$.
By \eqref{eq:48}, $|D(F_m)|=|D|^{p^m}$, so
$|D(F_m)|^{-1/d}=|D|^{-p^m/(d_0p^m)}=|D|^{-1/d_0}$, a constant independent of $m$.
Hence $\Bcal_d\ge C d$ with $C=|D|^{-1/d_0}$ --- the point being that the
discriminant grows exactly fast enough for the bound to stay linear rather than
decay.

The original typesets the exponent as $|D|^{-1_/d_0}$; and the sentence after
\eqref{eq:49} cites ``(3.10) et (4.8)'', where the pair actually being combined is the
upper bound \eqref{eq:310} and the lower bound \eqref{eq:49}.
\end{expansion}

If $p$ does not divide $D_F$, then $F$ and $\Q(\zeta_p)$ are linearly disjoint and one
may choose $E_p^1$ linearly disjoint from $F$. Then $F_0=F\,E_p^1$ and the inequality
is valid for $d=48(p-1)p^n$, $n\ge1$. Of course the factor $(p-1)$ is not necessary if
one is willing to use Artin's conjecture or Dedekind's divisibility conjecture.

\begin{expansion}
\textbf{The point of \S4, which is easy to lose.} \Ledger{E27} Theorem 3 gives
$B_d>d/(2\pi e)$ by a two-line argument about the volume of a ball. \S4 shows that,
along an infinite sequence of degrees, the \emph{same} linear growth is a consequence
of number theory: a field whose zeta function has a real zero forces
$\Bcal_d\ge d|D|^{-1/d}$, and class field towers produce fields of unboundedly large
degree whose root discriminants stay bounded. So the linear growth in Theorem 3 is not
an artefact of the crude volume bound --- it is the truth, and it was implicit in
known facts about ramification before anyone asked the Fourier question.
\end{expansion}

% =====================================================================
\section*{\S5. Editorial remarks --- everything later than the paper}
% =====================================================================

Quarantined here; nothing above depends on any of it.

\begin{enumerate}[leftmargin=1.6em,itemsep=5pt]
\item \textbf{The Hermite bound is $3/2$, not $3$.} \Ledger{C2} \S2.1 carries out the
  computation the paper describes but does not perform. With the paper's own Fourier
  convention the unique (up to scale) self-dual combination of $H_0,H_2,H_4$ with
  $f(0)=0$ is $f(x)=x^2\bigl(x^2-\tfrac3{2\pi}\bigr)e^{-\pi x^2}$, whose last sign
  change is at $\sqrt{3/(2\pi)}$, giving $\pi A^2=3/2$. The printed $\pi A^2\le3$ is
  true but weaker by a factor 2, and $3/2$ is exactly the bound $(2.5)$ obtained later
  from the $g_a$ family. So the two routes of \S2 --- Hermite, and the Gaussian
  family --- give the \emph{same} bound, which the paper does not remark on. We have
  not resolved whether the printed $3$ is a slip or reflects a different
  normalisation of $H_4$.

\item \textbf{Theorem 2 is weaker than its own method.} \Ledger{C21} The correction
  argument of \S2.3, optimised numerically, gives $X''=1{,}1512$ and $A\le0{,}6053$
  with $a_0=2$, and $X''\approx1{,}1210$, $A\le0{,}5973$ optimising $a_0$ as well;
  the paper states $X''=1{,}25\ldots$, $A\le0{,}63\ldots$ and Theorem 2 rounds to
  $A\le0{,}64$. The authors flag their computation as ``tr\`es approch\'e'', so this
  is a sharpening of an acknowledged estimate, not a correction. We state Theorem 2
  as they do.

\item \textbf{Where the true minimiser sits.} \Ledger{C19} $a=2{,}08137$ minimises the
  large-$a$ surrogate $\log(1+a)/(1-a^{-2})$; the true minimiser of $X_a$ is
  $a\approx1{,}89$, with $X\approx1{,}45200$ against $1{,}45341$ at $a=2$ and
  $1{,}45623$ at $a=2{,}08137$. All three are worse than the corrected value of
  item 2, which is the point of \S2.3.

\item \textbf{What the paper leaves open.} Two things, both explicit:
  \begin{itemize}[leftmargin=1.4em]
  \item The exact value of $B_1$. The paper brackets it as
    $0{,}1687\le B_1\le0{,}41$ and determines nothing further. The gap is a factor
    of about $2{,}4$.
  \item Whether $X(d)=\tfrac d2+1$ is optimal for $d\ge2$, i.e.\ whether
    $\Bcal_d=(d+2)/2\pi$. The paper says twice that this ``is possible'', having shown
    that the family $g_a$ cannot beat it for $d>2$ (via $Q_4(c+1)<0$) and for $d=2$
    (via $p_6<0$) --- in contrast with $d=1$, where it can.
  \end{itemize}
  Both are recorded on the project's Open problems page with what we could verify of
  their status.

\item \textbf{The linear growth constant is now known, and both of their bounds were
  within a factor of two.} Theorem 3 and Theorem 4 bracket $B_d$ between $d/2\pi e$ and
  $(d+2)/2\pi$, a factor $e$ apart. In August 2026 --- sixteen years later --- Chapter 1
  of OpenAI's \emph{Ten Advances in Mathematics and Theoretical Computer Science}
  states, as its Theorem 1.2, that
  \[\lim_{d\to\infty}\frac{A_+(d)}{\sqrt d}=\frac1\pi,\qquad\text{whence}\qquad
    \frac{B_d}{d}\longrightarrow\frac1{\pi^2}=0{,}101321\ldots\]
  where $A_+(d)$ is exactly the $A$ of \eqref{eq:34a}. The value sits inside their
  bracket, a factor $\pi/2=1{,}5708$ below their upper constant $1/2\pi=0{,}159155$ and
  a factor $2e/\pi=1{,}7305$ above their lower constant $1/2\pi e=0{,}058550$. So the
  crude volume bound of Theorem 3 and the Gaussian family of Theorem 4 were each wrong
  by less than a factor of two, and the truth lies between them, nearer the upper end.

  Two things must be said about that reference. Its results are stated to have been
  produced by an internal OpenAI model, and the manuscript is self-published rather
  than refereed. And we have read its statements, not its proof --- so this remark is
  reported, not verified, exactly like the imports of \S4. The same source's Theorem 1.2
  also settles the question of item 4 in the negative for large $d$: since
  $1/\pi^2<1/2\pi$, the value $X(d)=\tfrac d2+1$ is \emph{not} asymptotically optimal.
  Cohn--Gon\c calves had already shown it is not optimal at $d=12$, where the sharp
  value is $B_{12}=2$ against their $(12+2)/2\pi=2{,}2282$.

  For $d=1$ none of this helps: the exact value of $B_1$ is open, and
  Cohn--Gon\c calves record that no sharp value is known or even conjectured in any
  dimension but 12.

\item \textbf{The source.} The authors' arXiv LaTeX does not build on a 2026
  distribution: \verb|[francais]| must become \verb|[french]|, and two stray
  \texttt{0xCA} bytes (Mac Roman non-breaking spaces, at lines 58 and 400) must be
  replaced. With those two changes it compiles clean to exactly 19 pages, matching the
  arXiv PDF. The file is ISO-8859, not UTF-8. \Ledger{F2, F3}
\end{enumerate}

\end{document}
